
\documentclass{article}
\usepackage{colortbl}
\usepackage{makecell}
\usepackage{multirow}
\usepackage{supertabular}

\begin{document}

\newcounter{utterance}

\centering \large Interaction Transcript for game `eqbench', experiment `eq\_bench', episode 115 with reflexion{-}llama{-}3.1{-}8b{-}instruct.
\vspace{24pt}

{ \footnotesize  \setcounter{utterance}{1}
\setlength{\tabcolsep}{0pt}
\begin{supertabular}{c@{$\;$}|p{.15\linewidth}@{}p{.15\linewidth}p{.15\linewidth}p{.15\linewidth}p{.15\linewidth}p{.15\linewidth}}
    \# & \multicolumn{2}{c}{Player} && \multicolumn{2}{c}{Game Master} \\
    \hline

    \theutterance \stepcounter{utterance}  
    & & & \multicolumn{4}{p{0.6\linewidth}}{
        \cellcolor[rgb]{0.9,0.9,0.9}{
            \makecell[{{p{\linewidth}}}]{
                \texttt{\tiny{[P1$\langle$GM]}}
                \texttt{You have attempted this same task before and did not succeed. These are the plans you wrote after those attempts.} \\
\texttt{Trial \#1: New plan: } \\
\\ 
\texttt{Upon re{-}reading the dialogue, I realized that I focused too much on the surface{-}level interactions between Ben and Esther, without considering the underlying dynamics of their relationship and the context of the game they're playing. I should have paid more attention to the subtle cues and power dynamics at play.} \\
\\ 
\texttt{To improve my analysis, I will focus on the following key aspects:} \\
\\ 
\texttt{1. **Esther's character**: I should have delved deeper into Esther's personality, her motivations, and her emotional state. Her responses suggest a confident and assertive individual who is not easily intimidated. I will consider her past experiences and how they might influence her behavior in this situation.} \\
\texttt{2. **Ben's tactics**: I should have examined Ben's strategies more closely, particularly his attempts to provoke Esther and undermine her confidence. His comments about her "icy throne" and "lonely" life are likely intended to get under her skin and make her more emotional. I will consider how Ben's tactics might be affecting Esther's emotional state.} \\
\texttt{3. **The game context**: I should have taken into account the game of poker they're playing and how it might be influencing their interactions. The game is a high{-}stakes environment where players need to manage their emotions and make calculated decisions. I will consider how the game's dynamics might be affecting Ben and Esther's emotional responses.} \\
\\ 
\texttt{With these aspects in mind, I will re{-}evaluate the dialogue and provide a revised analysis.} \\
\\ 
\texttt{Trial \#2: New plan: } \\
\\ 
\texttt{Upon re{-}reading the dialogue, I realized that I focused too much on the surface{-}level interactions between Ben and Esther, without considering the underlying dynamics of their relationship and the context of the game they're playing. I should have paid more attention to the subtle cues and power dynamics at play.} \\
\\ 
\texttt{To improve my analysis, I will focus on the following key aspects:} \\
\\ 
\texttt{1. **Esther's character**: I should have delved deeper into Esther's personality, her motivations, and her emotional state. Her responses suggest a confident and assertive individual who is not easily intimidated. I will consider her past experiences and how they might influence her behavior in this situation.} \\
\texttt{2. **Ben's tactics**: I should have examined Ben's strategies more closely, particularly his attempts to provoke Esther and undermine her confidence. His comments about her "icy throne" and "lonely" life are likely intended to get under her skin and make her more emotional. I will consider how Ben's tactics might be affecting Esther's emotional state.} \\
\texttt{3. **The game context**: I should have taken into account the game of poker they're playing and how it might be influencing their interactions. The game is a high{-}stakes environment where players need to manage their emotions and make calculated decisions. I will consider how the game's dynamics might be affecting Ben and Esther's emotional responses.} \\
\\ 
\texttt{To improve my analysis, I will also consider the following:} \\
\\ 
\texttt{* I will look for subtle cues in Esther's language and tone that might indicate her emotional state, such as her use of sarcasm or her body language.} \\
\texttt{* I will examine Ben's behavior more closely, particularly his attempts to manipulate Esther's emotions and make her more emotional.} \\
\texttt{* I will consider the power dynamics at play in the game, including Esther's confidence and experience, and how they might be influencing Ben's behavior.} \\
\\ 
\texttt{With these aspects in mind, I will re{-}evaluate the dialogue and provide a revised analysis.} \\
\\ 
\\ 
\\ 
\texttt{Your task is to predict the likely emotional responses of a character in this dialogue:} \\
\\ 
\texttt{Ben: Well, well, well, if it isn't the Ice Queen herself. I've heard tales of your cold{-}blooded poker face.} \\
\texttt{Esther: And you must be Benny the Bluff. I've heard you're quite the charmer. } \\
\texttt{Ben: Guilty as charged, sweetheart. But don't let that fool you. I've got a few tricks up my sleeve.} \\
\texttt{Esther: I'm sure you do. But remember, Ben, this isn't a magic show. It's a game of skill. } \\
\texttt{Ben: And a little bit of luck, wouldn't you agree?} \\
\texttt{Esther: Luck is for those who can't rely on their abilities.} \\
\texttt{Ben: Spoken like a true Ice Queen. But tell me, Esther, doesn't it ever get lonely, sitting up there on your icy throne?} \\
\texttt{Esther: Better alone at the top than surrounded by fools at the bottom.} \\
\texttt{Ben: Ouch! That's a sharp tongue you've got there. } \\
\texttt{Esther: It's served me well so far.} \\
\texttt{Ben: I can see that. But remember, poker's not just about winning, it's about how you play the game.} \\
\texttt{Esther: And yet, it seems I'm the one raking in the chips.} \\
\texttt{Ben: For now, sweetheart, for now. But the night's still young.} \\
\texttt{Esther: Indeed, it is. Let's see if your bluffs can keep up with my bets. } \\
\texttt{Ben: Oh, they will, Esther. They always do.} \\
\texttt{Esther: We'll see about that, Benny. We'll see.} \\
\texttt{{[}End dialogue{]}} \\
\\ 
\texttt{At the end of this dialogue, Ben would feel...} \\
\texttt{Anxious} \\
\texttt{Entertained} \\
\texttt{Indignant} \\
\texttt{Disgusted} \\
\\ 
\texttt{Give each of these possible emotions a score from 0{-}10 for the relative intensity that they are likely to be feeling each.} \\
\\ 
\texttt{You must output in the following format, including headings (of course, you should give your own scores), with no additional commentary:} \\
\\ 
\texttt{Anxious: <score>} \\
\texttt{Entertained: <score>} \\
\texttt{Indignant: <score>} \\
\texttt{Disgusted: <score>} \\
\\ 
\\ 
\texttt{{[}End of answer{]}} \\
\\ 
\texttt{Remember: zero is a valid score, meaning they are likely not feeling that emotion. You must score at least one emotion > 0.} \\
\\ 
\texttt{Your answer:} \\
            }
        }
    }
     \\ \\

    \theutterance \stepcounter{utterance}  
    & \multicolumn{4}{p{0.6\linewidth}}{
        \cellcolor[rgb]{0.561,0.737,0.561}{
            \makecell[{{p{\linewidth}}}]{
                \texttt{\tiny{[P1$\rangle$GM]}}
                \texttt{Anxious: 2} \\
\texttt{Entertained: 8} \\
\texttt{Indignant: 6} \\
\texttt{Disgusted: 0} \\
            }
        }
    }
    & & \\ \\

    \theutterance \stepcounter{utterance}  
    & & & \multicolumn{2}{p{0.3\linewidth}}{
        \cellcolor[rgb]{0.9,0.9,0.9}{
            \makecell[{{p{\linewidth}}}]{
                \texttt{\tiny{[GM$|$GM]}}
                \texttt{Anxious: 2} \\
\texttt{Entertained: 8} \\
\texttt{Indignant: 6} \\
\texttt{Disgusted: 0} \\
            }
        }
    }
    & & \\ \\

    \theutterance \stepcounter{utterance}  
    & & & \multicolumn{2}{p{0.3\linewidth}}{
        \cellcolor[rgb]{0.9,0.9,0.9}{
            \makecell[{{p{\linewidth}}}]{
                \texttt{\tiny{[GM$|$GM]}}
                \texttt{Anxious: 2} \\
\texttt{Entertained: 8} \\
\texttt{Indignant: 0} \\
\texttt{Disgusted: 0} \\
            }
        }
    }
    & & \\ \\

    \theutterance \stepcounter{utterance}  
    & & & \multicolumn{2}{p{0.3\linewidth}}{
        \cellcolor[rgb]{0.9,0.9,0.9}{
            \makecell[{{p{\linewidth}}}]{
                \texttt{\tiny{[GM$|$GM]}}
                \texttt{game\_result = LOSE} \\
            }
        }
    }
    & & \\ \\

    \theutterance \stepcounter{utterance}  
    & & \multicolumn{4}{p{0.6\linewidth}}{
        \cellcolor[rgb]{0.95,0.95,0.95}{
            \makecell[{{p{\linewidth}}}]{
                \texttt{\tiny{[P1$\leftrightarrow$P]}}
                \texttt{Post{-}episode Reflexion (not sent to GM):} \\
\texttt{New plan: } \\
\\ 
\texttt{Upon re{-}reading the dialogue, I realized that I focused too much on the surface{-}level interactions between Ben and Esther, without considering the underlying dynamics of their relationship and the context of the game they're playing. I should have paid more attention to the subtle cues and power dynamics at play.} \\
\\ 
\texttt{To improve my analysis, I will focus on the following key aspects:} \\
\\ 
\texttt{1. **Esther's character**: I should have delved deeper into Esther's personality, her motivations, and her emotional state. Her responses suggest a confident and assertive individual who is not easily intimidated. I will consider her past experiences and how they might influence her behavior in this situation.} \\
\texttt{2. **Ben's tactics**: I should have examined Ben's strategies more closely, particularly his attempts to provoke Esther and undermine her confidence. His comments about her "icy throne" and "lonely" life are likely intended to get under her skin and make her more emotional. I will consider how Ben's tactics might be affecting Esther's emotional state.} \\
\texttt{3. **The game context**: I should have taken into account the game of poker they're playing and how it might be influencing their interactions. The game is a high{-}stakes environment where players need to manage their emotions and make calculated decisions. I will consider how the game's dynamics might be affecting Ben and Esther's emotional responses.} \\
\\ 
\texttt{To improve my analysis, I will also consider the following:} \\
\\ 
\texttt{* I will look for subtle cues in Esther's language and tone that might indicate her emotional state, such as her use of sarcasm or her body language.} \\
\texttt{* I will examine Ben's behavior more closely, particularly his attempts to manipulate Esther's emotions and make her more emotional.} \\
\texttt{* I will consider the power dynamics at play in the game, including Esther's confidence and experience, and how they might be influencing Ben's behavior.} \\
\\ 
\texttt{To provide a more accurate analysis, I will re{-}evaluate the dialogue and consider the following:} \\
\\ 
\texttt{* Esther's confidence and assertiveness suggest that she is not easily intimidated, but Ben's comments might be trying to get under her skin and make her more emotional.} \\
\texttt{* Ben's tactics, such as making personal comments and trying to provoke Esther, are likely intended to undermine her confidence and make her more emotional.} \\
\texttt{* The game context is a high{-}stakes environment where players need to manage their emotions and make calculated decisions, which might be influencing Ben and Esther's emotional responses.} \\
\\ 
\texttt{With these aspects in mind, I will re{-}evaluate the dialogue and provide a revised analysis.} \\
\\ 
\texttt{Anxious: 0} \\
\texttt{Entertained: 9} \\
\texttt{Indignant: 7} \\
\texttt{Disgusted: 0} \\
            }
        }
    }
    & \\ \\

\end{supertabular}
}

\end{document}
